\begin{figure}[ht]
\centering
\begin{tikzpicture}[x=1cm, y=1cm, font=\small]

% Term magnitudes of the two arctan series, log scale, vs term index k.
% arctan(1/5): term_k = (1/5)^(2k+1)/(2k+1). log10 values:
%  k=0: 1/5=0.2 -> -0.70; k=1: (0.2)^3/3=2.67e-3 -> -2.57; k=2: (0.2)^5/5=6.4e-5 -> -4.19;
%  k=3: (0.2)^7/7=1.83e-6 -> -5.74; k=4: 5.69e-8 -> -7.24; k=5: 1.86e-9 -> -8.73;
%  k=6: 6.3e-11 -> -10.2; k=7: 2.17e-12 -> -11.66
% arctan(1/239): term_k = (1/239)^(2k+1)/(2k+1):
%  k=0: 4.18e-3 -> -2.38; k=1: 2.44e-8 -> -7.61; k=2: 8.6e-14 -> -13.07
% y scale: log10 from 0 (top) to -14 (bottom). y_cm = (val+14)*0.55 ; map.

% Axes (y from -14 bottom to 0 top, total 14 decades; use y_cm = val (negative) shifted)
\draw[->, thick] (-0.3, -14) -- (8.0, -14) node[right] {term index $k$};
\draw[->, thick] (-0.3, -14.3) -- (-0.3, 0.8) node[above] {term magnitude (log scale)};

% y ticks every 2 decades
\foreach \p in {0,-2,-4,-6,-8,-10,-12,-14} {
  \draw (-0.45, \p) -- (-0.15, \p) node[left, xshift=-0.1cm] {$10^{\p}$};
}
% x ticks at k=0..7 mapped x_cm = k*1.0
\foreach \k in {0,1,2,3,4,5,6,7} {
  \draw (\k, -14.15) -- (\k, -13.85) node[below, yshift=-0.1cm] {\k};
}

% per-series tolerance line at 5e-12 -> log10 = -11.30
\draw[dashed, gray] (-0.3, -11.30) -- (7.6, -11.30)
  node[right, gray, yshift=0.18cm] {budget $5\times10^{-12}$};

% arctan(1/5) series
\draw[thick, engblue, mark=*, mark size=1pt] plot coordinates {
  (0,-0.70) (1,-2.57) (2,-4.19) (3,-5.74) (4,-7.24) (5,-8.73)
  (6,-10.2) (7,-11.66) };
\node[engblue, anchor=west] at (5.1, -8.0) {$\arctan\tfrac{1}{5}$};

% arctan(1/239) series
\draw[thick, engred, mark=square*, mark size=1pt] plot coordinates {
  (0,-2.38) (1,-7.61) (2,-13.07) };
\node[engred, anchor=west] at (2.1, -13.07) {$\arctan\tfrac{1}{239}$};

\end{tikzpicture}
\caption[Truncation budget for the Machin computation of $\pi$]{Term
magnitudes of the two arctangent series in Machin's formula $\pi/4 = 4
\arctan\tfrac{1}{5} - \arctan\tfrac{1}{239}$, on a log scale against term
index. The dashed line is the per-series truncation budget. Each series is
cut once the next-term bound drops below the line: at index seven for the
slowly converging $1/5$ series, at index one for the fast $1/239$ series
(whose index-two term already sits far below budget). The asymmetry is the
geometric convergence rate at work, $1/25$ per step against $1/57121$ per
step. The figure is the case study's error budget made visual.}
\label{fig:vol02:ch04:machin-budget}
\end{figure}
